% !TEX program = xelatex
% Lernplatt - by Can Siebert
% Kurs 22 (Bo 143) - Tag 2 - Aufgabenheft
% Plattform-Implementierung, KPIs & strategische Weiterentwicklung
%
% Self-contained xelatex source. Kein biblatex, keine externen Bilder.

\documentclass[11pt,a4paper,oneside]{article}

% --- Encoding & Sprache --------------------------------------------------
\usepackage{polyglossia}
\setdefaultlanguage{german}

% --- Geometrie -----------------------------------------------------------
\usepackage[a4paper,margin=2.5cm]{geometry}

% --- Fonts ---------------------------------------------------------------
\usepackage{fontspec}
\defaultfontfeatures{Ligatures=TeX}

\IfFontExistsTF{Charter}{%
  \setmainfont{Charter}%
}{%
  \IfFontExistsTF{Noto Serif}{%
    \setmainfont{Noto Serif}%
  }{%
    \setmainfont{Liberation Serif}%
  }%
}

\IfFontExistsTF{Source Sans 3}{%
  \setsansfont{Source Sans 3}%
}{%
  \IfFontExistsTF{Source Sans Pro}{%
    \setsansfont{Source Sans Pro}%
  }{%
    \setsansfont{Liberation Sans}%
  }%
}

\newcommand{\caveatfontfallback}{\itshape}
\IfFontExistsTF{Caveat}{%
  \newfontfamily\caveatfont{Caveat}[Scale=1.15]%
}{%
  \newcommand\caveatfont{\caveatfontfallback}%
}

\setmonofont{Liberation Mono}

% --- Mikro-Typografie ----------------------------------------------------
\usepackage{microtype}
\usepackage{eurosym}
\linespread{1.15}

% --- Farben & Boxen ------------------------------------------------------
\usepackage{xcolor}
\definecolor{lpInk}{RGB}{20,28,38}
\definecolor{lpAccent}{RGB}{22,90,140}
\definecolor{lpAccentLight}{RGB}{210,228,240}
\definecolor{lpAmber}{RGB}{222,138,30}
\definecolor{lpAmberLight}{RGB}{255,243,217}
\definecolor{lpAmberBorder}{RGB}{196,118,18}
\definecolor{lpGray}{RGB}{96,104,114}
\definecolor{lpGrayLight}{RGB}{238,240,243}
\definecolor{lpGreen}{RGB}{34,120,72}
\definecolor{lpGreenLight}{RGB}{222,238,228}
\definecolor{lpGreenBorder}{RGB}{24,96,58}

% Niveau-Farben
\definecolor{lpL1}{RGB}{34,120,72}      % L1 = gruen (reproduktion)
\definecolor{lpL2}{RGB}{22,90,140}      % L2 = blau (anwendung)
\definecolor{lpL3}{RGB}{170,42,52}      % L3 = rot (management)

\usepackage[most]{tcolorbox}
\tcbuselibrary{breakable,skins}

% Info-Box (blau-weiss) fuer Hinweise/Aufgabenstellung-Rahmen
\newtcolorbox{infobox}[1][Hinweis]{%
  enhanced, breakable,
  colback=lpAccentLight,
  colframe=lpAccent,
  boxrule=0.6pt,
  arc=2pt,
  left=10pt,right=10pt,top=8pt,bottom=8pt,
  fonttitle=\sffamily\bfseries\color{white},
  coltitle=white,
  title={#1},
  attach boxed title to top left={xshift=8pt,yshift=-8pt},
  boxed title style={size=small, colback=lpAccent, colframe=lpAccent, boxrule=0.4pt, arc=1pt},
  top=14pt
}

% Antwort-Bereich: leerer Kasten mit Punktlinien
\newtcolorbox{antwortbox}[1][Ihre Antwort]{%
  enhanced, breakable,
  colback=white,
  colframe=lpGray,
  boxrule=0.4pt,
  arc=2pt,
  left=10pt,right=10pt,top=8pt,bottom=8pt,
  fonttitle=\sffamily\bfseries\color{white},
  coltitle=white,
  title={#1},
  attach boxed title to top left={xshift=8pt,yshift=-8pt},
  boxed title style={size=small, colback=lpGray, colframe=lpGray, boxrule=0.4pt, arc=1pt},
  top=14pt
}

% --- Listen --------------------------------------------------------------
\usepackage{enumitem}
\setlist{itemsep=2pt, topsep=4pt, parsep=0pt}
\setlist[itemize,1]{label={\textbullet},leftmargin=1.6em}
\setlist[itemize,2]{label={\textendash},leftmargin=1.4em}
\setlist[enumerate,1]{label=\arabic*., leftmargin=1.8em}

% --- Tabellen ------------------------------------------------------------
\usepackage{tabularx}
\usepackage{array}
\usepackage{booktabs}
\usepackage{longtable}

% --- Kopf-/Fusszeilen ----------------------------------------------------
\usepackage{fancyhdr}
\usepackage{lastpage}
\pagestyle{fancy}
\fancyhf{}
\renewcommand{\headrulewidth}{0.3pt}
\renewcommand{\footrulewidth}{0pt}
\fancyhead[L]{\footnotesize\sffamily\color{lpGray}Aufgabenheft \textperiodcentered{} Kurs 22 \textperiodcentered{} Tag 2}
\fancyhead[R]{\footnotesize\sffamily\color{lpGray}Plattform-Implementierung \& KPIs}
\fancyfoot[L]{\footnotesize\sffamily\color{lpGray}\textcopyright{} Can Siebert}
\fancyfoot[R]{\footnotesize\sffamily\color{lpGray}\thepage{} / \pageref{LastPage}}

\fancypagestyle{plain}{%
  \fancyhf{}%
  \renewcommand{\headrulewidth}{0pt}%
  \fancyfoot[C]{\footnotesize\sffamily\color{lpGray}\thepage{} / \pageref{LastPage}}%
}

% --- Section-Styling -----------------------------------------------------
\usepackage[explicit]{titlesec}

\titleformat{\section}[block]
  {\sffamily\Large\bfseries\color{lpAccent}}
  {\thesection.}{0.6em}{#1}
  [{\vspace{2pt}\color{lpAccent}\titlerule[0.6pt]}]

\titleformat{\subsection}[block]
  {\sffamily\large\bfseries\color{lpInk}}
  {\thesubsection}{0.6em}{#1}

\titlespacing*{\section}{0pt}{14pt}{8pt}
\titlespacing*{\subsection}{0pt}{10pt}{4pt}

% --- Hyperlinks ----------------------------------------------------------
\usepackage[hidelinks,unicode]{hyperref}

% --- Helfer --------------------------------------------------------------
\newcommand{\akzent}[1]{{\caveatfont\color{lpAmber}#1}}
\newcommand{\fachbegriff}[1]{\textbf{#1}}

% Niveau-Badge: \niveau{L1}, \niveau{L2}, \niveau{L3}
\newcommand{\niveaubadge}[2]{%
  \tcbox[on line, colback=#1, colframe=#1, coltext=white,
         boxrule=0pt, arc=2pt, left=5pt,right=5pt,top=1pt,bottom=1pt,
         nobeforeafter]{\sffamily\bfseries\footnotesize #2}%
}
\newcommand{\nivLi}{\niveaubadge{lpL1}{L1\,\textbullet\,Wissen}}
\newcommand{\nivLii}{\niveaubadge{lpL2}{L2\,\textbullet\,Anwendung}}
\newcommand{\nivLiii}{\niveaubadge{lpL3}{L3\,\textbullet\,Management}}

% Zeit-Hinweis
\newcommand{\zeit}[1]{%
  \tcbox[on line, colback=lpGrayLight, colframe=lpGrayLight, coltext=lpInk,
         boxrule=0pt, arc=2pt, left=5pt,right=5pt,top=1pt,bottom=1pt,
         nobeforeafter]{\sffamily\footnotesize\textbf{$\sim$\,#1}}%
}

% Aufgaben-Kopf: Nr + Titel + Badges + Zeit
\newcommand{\aufgabenkopf}[4]{%
  \par\vspace{6pt}
  \noindent{\sffamily\Large\bfseries\color{lpAccent}Aufgabe #1 \textendash{} #2}\par
  \vspace{2pt}
  \noindent{\color{lpAccent}\rule{\linewidth}{0.6pt}}\par
  \vspace{4pt}
  \noindent #3 \hfill \zeit{#4}\par
  \vspace{6pt}
}

% YouTube-Suchquery (kompakt rendern)
\newcommand{\ytquery}[1]{%
  \par\vspace{4pt}
  \noindent{\footnotesize\sffamily\color{lpGray}\textbf{Lernvideo zum Thema:}\ %
    \href{https://studyflix.de/search?query=#1}{\textcolor{lpAccent}{Studyflix-Treffer}}\ ·\ %
    \href{https://www.youtube.com/results?search\_query=#1}{\textcolor{lpAccent}{YouTube-Treffer}}\\%
    \textit{\textcolor{lpGray}{Stichworte: #1}}}\par
}

% Linierter Antwort-Bereich mit gepunkteten Linien
\newcommand{\antwortlinien}[1]{%
  \par\vspace{6pt}
  \foreach \x in {1,...,#1}{%
    \noindent\makebox[\linewidth]{\dotfill}\\[6pt]
  }
}
\usepackage{pgffor}

% =========================================================================
\begin{document}

% --- COVER ---------------------------------------------------------------
\begin{titlepage}
  \pagestyle{empty}
  \vspace*{1.5cm}
  \noindent{\sffamily\footnotesize\color{lpGray}LERNPLATT \textperiodcentered{} BY CAN SIEBERT}\\[2pt]
  \noindent{\color{lpAccent}\rule{\linewidth}{1.2pt}}\\[4pt]
  \noindent{\sffamily\footnotesize\color{lpGray}AUFGABENHEFT \textperiodcentered{} MODUL 1 \textperiodcentered{} TAG 2 \textperiodcentered{} KURS 22 (Bo 143) \textperiodcentered{} 29.\,Mai 2026}

  \vspace{3cm}
  {\sffamily\fontsize{30}{34}\selectfont\bfseries\color{lpInk}Aufgabenheft\\Tag 2\par}

  \vspace{0.7cm}
  {\sffamily\large\color{lpGray}Plattform-Implementierung, KPIs\\\& strategische Weiterentwicklung\\\textendash{} elf Aufgaben auf drei Niveaus.\par}

  \vspace{1.5cm}
  {\caveatfont\LARGE\color{lpAmber}11 Aufgaben \textperiodcentered{} L1 \textperiodcentered{} L2 \textperiodcentered{} L3}\par

  \vfill

  \noindent\begin{minipage}{0.55\linewidth}
    {\sffamily\footnotesize\color{lpGray}Niveau-Mix:}\\
    {\sffamily\small\color{lpInk}3\,\textsf{L1} \textperiodcentered{} 5\,\textsf{L2} \textperiodcentered{} 3\,\textsf{L3}}\\[10pt]
    {\sffamily\footnotesize\color{lpGray}Bearbeitung:}\\
    {\sffamily\small\color{lpInk}Selbstbearbeitung im Lernpfad}
  \end{minipage}\hfill
  \begin{minipage}{0.4\linewidth}
    \raggedleft
    {\sffamily\footnotesize\color{lpGray}Dozent:}\\
    {\sffamily\small\color{lpInk}Can Siebert}\\[10pt]
    {\sffamily\footnotesize\color{lpGray}Begleitend:}\\
    {\sffamily\small\color{lpInk}Lehrskript \& Lösungsheft}
  \end{minipage}

  \vspace{0.5cm}
  \noindent{\color{lpAccent}\rule{\linewidth}{0.6pt}}
\end{titlepage}

% --- ANLEITUNG -----------------------------------------------------------
\section*{So arbeiten Sie mit diesem Heft}
\thispagestyle{plain}

\noindent Dieses Aufgabenheft enthält elf Aufgaben zu Tag 2 \textendash{} Plattform-Implementierung, KPI-Steuerung und strategische Weiterentwicklung. Tag 2 baut auf Tag 1 auf: dort wurde die \emph{Auswahl} einer Plattformstrategie geübt, heute geht es um die \emph{Umsetzung, Steuerung und Optimierung}.

\begin{itemize}
  \item \nivLi{} \textendash{} \fachbegriff{Wissen.} Implementierungsphasen, KPI-Logik und Zukunftstrends reproduzieren. 5\,--\,10 Minuten pro Aufgabe.
  \item \nivLii{} \textendash{} \fachbegriff{Anwendung.} Transfer auf SmartOffice24, BauTech Connect und konkrete Plattform-Cases. 15\,--\,30 Minuten.
  \item \nivLiii{} \textendash{} \fachbegriff{Management.} Entscheidungsarchitektur, Verantwortung, Plattformsteuerung bis 2030. 30\,--\,45 Minuten.
\end{itemize}

\begin{infobox}[Hinweis]
Die Musterlösungen finden Sie im separaten \emph{Lösungsheft Tag 2}. Versuchen Sie jede Aufgabe zuerst eigenständig \textendash{} der Mehrwert entsteht in der eigenen Argumentation, nicht im Abschreiben.
\end{infobox}

\clearpage

% =========================================================================
% L1 AUFGABEN
% =========================================================================
\section*{Niveau L1 \textendash{} Wissen \& Reproduktion}
\addcontentsline{toc}{section}{L1 -- Wissen}

% --- AUFGABE 1 -----------------------------------------------------------
% LERNPLATT_HILFSVIDEOS
\section*{Hilfsvideos zu diesem Tag}
\addcontentsline{toc}{section}{Hilfsvideos zu diesem Tag}
\thispagestyle{plain}

\noindent Die folgenden YouTube-Erkl\"arvideos vermitteln die Inhalte, die Sie zur Bearbeitung der Aufgaben ben\"otigen. Klicken Sie auf den Titel, um das Video direkt zu \"offnen; die vollst\"andige URL steht in Klammern.

\begin{description}[style=nextline,leftmargin=18pt,labelindent=0pt,itemsep=8pt]
  \item[1.~Lean Startup — Build, Measure, Learn als Pilotierungs-Logik] \href{https://youtu.be/vJItJyGp4VI}{\textcolor{lpAccent}{https://youtu.be/vJItJyGp4VI}}\\[2pt]
  {\footnotesize\color{lpGray}(URL: \nolinkurl{https://youtu.be/vJItJyGp4VI})}
  \item[2.~Conversion Rate Optimierung im E-Commerce — die Praxis] \href{https://youtu.be/JgA_nNf_XcY}{\textcolor{lpAccent}{\nolinkurl{https://youtu.be/JgA_nNf_XcY}}}\\[2pt]
  {\footnotesize\color{lpGray}(URL: \nolinkurl{https://youtu.be/JgA_nNf_XcY})}
  \item[3.~Customer Lifetime Value — was ein Kunde wirklich wert ist] \href{https://youtu.be/cmhmiDWztzU}{\textcolor{lpAccent}{https://youtu.be/cmhmiDWztzU}}\\[2pt]
  {\footnotesize\color{lpGray}(URL: \nolinkurl{https://youtu.be/cmhmiDWztzU})}
  \item[4.~KI im B2B-Vertrieb — vom Trend zur konkreten Anwendung] \href{https://youtu.be/PShGekYpcwI}{\textcolor{lpAccent}{https://youtu.be/PShGekYpcwI}}\\[2pt]
  {\footnotesize\color{lpGray}(URL: \nolinkurl{https://youtu.be/PShGekYpcwI})}
\end{description}
\vspace{0.6em}
\noindent{\footnotesize\color{lpGray}\textit{Tipp:} \"Offnen Sie das Video, lassen Sie es im Hintergrund laufen und arbeiten Sie parallel an der Aufgabe.}

\clearpage

\aufgabenkopf{1}{Auswahl \textendash{} Implementierung \textendash{} Optimierung}{\nivLi}{6\,min}

\noindent Drei Begriffe werden im Plattformkontext häufig vermischt: \emph{Plattformauswahl}, \emph{Plattformimplementierung} und \emph{Plattformoptimierung}.

\medskip
\noindent\textbf{Aufgabe:}

\begin{enumerate}
  \item Definieren Sie jeden der drei Begriffe in \emph{einem} Satz.
  \item Schreiben Sie zu jedem Begriff einen typischen Stolperstein, an dem Mittelständler häufig scheitern.
  \item Ergänzen Sie einen Schlusssatz: Warum ist es entscheidend, diese drei Phasen sauber zu trennen?
\end{enumerate}

\ytquery{Plattformauswahl Plattformimplementierung Plattformoptimierung Unterschied}

\antwortlinien{8}

\clearpage

% --- AUFGABE 2 -----------------------------------------------------------
\aufgabenkopf{2}{Sechs Phasen einer Plattformimplementierung}{\nivLi}{8\,min}

\noindent Eine geordnete Plattformimplementierung umfasst typischerweise sechs Phasen: \emph{Zielbild \textperiodcentered{} Pilotierung \textperiodcentered{} Rollout \textperiodcentered{} Betrieb \textperiodcentered{} Skalierung \textperiodcentered{} Optimierung}.

\medskip
\noindent\textbf{Aufgabe:}

\begin{enumerate}
  \item Ordnen Sie die folgenden sechs Tätigkeiten der jeweils passenden Phase zu:
  \begin{itemize}
    \item Kennzahlen für Sichtbarkeit, Conversion und Wiederkauf auswerten und Anpassungen testen.
    \item Mit drei Pilotkunden und einem reduzierten Sortiment den Plattformprozess validieren.
    \item Die strategische Zielsetzung formulieren: \emph{Welchen Kundenproblem löst die Plattform und welche Rolle nimmt sie im Vertrieb ein?}
    \item Stabile Servicelevel sichern: Antwortzeiten, Reklamationen, Verfügbarkeit, CRM-Pflege.
    \item Auf weitere Regionen, Sortimente oder Marktseiten ausweiten.
    \item Vom Pilot in den breiten Marktstart überführen \textendash{} inklusive Schulungen und Rollenklärung.
  \end{itemize}
  \item Benennen Sie für \emph{jede} Phase eine typische Verantwortlichkeit (Vertrieb, Marketing, IT, Kundenservice, Geschäftsführung).
\end{enumerate}

\ytquery{Plattformimplementierung Phasen Zielbild Pilot Rollout Betrieb Skalierung}

\antwortlinien{10}

\clearpage

% --- AUFGABE 3 -----------------------------------------------------------
\aufgabenkopf{3}{Plattform-KPIs sauber unterscheiden}{\nivLi}{8\,min}

\noindent Im Plattformbetrieb werden oft zu viele Kennzahlen gleichzeitig betrachtet \textendash{} und am Ende steuert niemand mit ihnen. Üblich sind unter anderem: \emph{Reichweite \textperiodcentered{} Traffic \textperiodcentered{} Leads \textperiodcentered{} Conversion Rate \textperiodcentered{} Kundenakquisitionskosten (CAC) \textperiodcentered{} Wiederkaufrate \textperiodcentered{} Warenkorbwert \textperiodcentered{} Churn Rate \textperiodcentered{} Customer Lifetime Value (CLV) \textperiodcentered{} Plattformkosten \textperiodcentered{} Deckungsbeitrag}.

\medskip
\noindent\textbf{Aufgabe:}

\begin{enumerate}
  \item Ordnen Sie jede der elf Kennzahlen \emph{einer} dieser drei Kategorien zu:
  \begin{itemize}
    \item \emph{Sichtbarkeitskennzahl} (zeigt, ob die Plattform überhaupt gefunden wird)
    \item \emph{Conversion-/Aktivierungskennzahl} (zeigt, ob aus Sichtbarkeit Geschäft wird)
    \item \emph{Wirtschaftlichkeitskennzahl} (zeigt, ob das Geschäft profitabel und nachhaltig ist)
  \end{itemize}
  \item Benennen Sie für jede der drei Kategorien eine Kennzahl, die ein \emph{operativer Plattformmanager} täglich braucht \textendash{} und eine, die im \emph{Management-Board} berichtet werden muss.
\end{enumerate}

\ytquery{Plattform KPIs Conversion CAC CLV Wiederkaufrate Deckungsbeitrag}

\antwortlinien{12}

\clearpage

% =========================================================================
% L2 AUFGABEN
% =========================================================================
\section*{Niveau L2 \textendash{} Anwendung \& Transfer}
\addcontentsline{toc}{section}{L2 -- Anwendung}

% --- AUFGABE 4 -----------------------------------------------------------
\aufgabenkopf{4}{SmartOffice24 \textendash{} Implementierung in drei Phasen}{\nivLii}{25\,min}

\begin{infobox}[Fallbeschreibung]
\textbf{Unternehmen:} ``SmartOffice24'' verkauft hochwertige Bürotechnik an kleine und mittlere Unternehmen. Bisheriger Vertrieb: Außendienst und Telefon. Die Geschäftsleitung hat entschieden, künftig zusätzlich über einen eigenen B2B-Webshop und einen externen B2B-Marktplatz zu verkaufen.

\smallskip
\textbf{Gegeben:}
\begin{itemize}
  \item Der Webshop soll in 6 Monaten starten.
  \item Der Marktplatz soll bereits in 3 Monaten erste Verkäufe ermöglichen.
  \item Produktdaten sind unvollständig.
  \item Das Vertriebsteam ist skeptisch.
  \item Marketing hat wenig Erfahrung mit Plattformen.
  \item IT-Ressourcen sind begrenzt.
  \item Budget: 80.000\,\euro{} für die erste Umsetzungsphase.
\end{itemize}
\end{infobox}

\noindent\textbf{Arbeitsauftrag:}

\begin{enumerate}
  \item Gliedern Sie die Implementierung in \textbf{drei sinnvolle Phasen} mit klaren Meilensteinen über die ersten 6 Monate. Begründen Sie, warum diese Reihenfolge sinnvoll ist.
  \item Benennen Sie mindestens \textbf{fünf Aufgaben}, die \emph{vor} dem Plattformstart erledigt werden müssen.
  \item Ordnen Sie jede Aufgabe genau einem Bereich zu: \emph{Vertrieb \textperiodcentered{} Marketing \textperiodcentered{} IT \textperiodcentered{} Kundenservice \textperiodcentered{} Management}.
  \item Identifizieren Sie \textbf{drei Risiken}, die den Start gefährden, und je eine Gegenmaßnahme.
  \item Treffen Sie eine Entscheidung: Sollte der Marktplatz wirklich \emph{vor} dem Webshop live gehen \textendash{} oder ist die Reihenfolge falsch? Begründen Sie in 3\,--\,4 Sätzen.
\end{enumerate}

\ytquery{B2B Webshop Marktplatz Implementierung Mittelstand Phasen Risiken}

\antwortlinien{20}

\clearpage

% --- AUFGABE 5 -----------------------------------------------------------
\aufgabenkopf{5}{SmartOffice24 \textendash{} Plattformkennzahlen interpretieren}{\nivLii}{30\,min}

\begin{infobox}[Kennzahlen nach 3 Monaten Plattformbetrieb]
\textbf{B2B-Webshop:}
\begin{itemize}
  \item 12.000 Besucher \textperiodcentered{} 480 registrierte Nutzer \textperiodcentered{} 120 Bestellungen
  \item Durchschnittlicher Bestellwert: 420\,\euro
  \item Conversion Rate: 1,0\,\%
  \item Auffällig: viele Abbrüche im Checkout
\end{itemize}

\textbf{Externer B2B-Marktplatz:}
\begin{itemize}
  \item 35.000 Produktaufrufe \textperiodcentered{} 300 Bestellungen
  \item Durchschnittlicher Bestellwert: 190\,\euro
  \item Hoher Preisdruck
  \item Geringer Zugriff auf Kundendaten
\end{itemize}

\textbf{Rahmenbedingungen:}
\begin{itemize}
  \item Geschäftsführung erwartet schnelles Wachstum.
  \item Vertrieb möchte hochwertige Kundenbeziehungen schützen.
  \item Marketing fordert mehr Budget für Anzeigen.
  \item IT weist auf technische Begrenzungen im Webshop hin.
\end{itemize}
\end{infobox}

\noindent\textbf{Arbeitsauftrag:}

\begin{enumerate}
  \item Beschreiben Sie \textbf{Stärken und Schwächen} beider Plattformkanäle (jeweils 2\,--\,3 Punkte).
  \item Benennen Sie \textbf{drei mögliche Ursachen} für die niedrige Conversion Rate im Webshop.
  \item Entwickeln Sie \textbf{drei Optimierungsmaßnahmen} für den Webshop \textendash{} priorisiert nach Wirkung und Umsetzbarkeit.
  \item Entwickeln Sie \textbf{zwei Maßnahmen} zur besseren Steuerung des Marktplatzgeschäfts (Sortiment, Marge, Daten).
  \item Entscheiden Sie, wohin das nächste Budget fließen soll \textendash{} Webshop, Marktplatz oder Marketing \textendash{} und begründen Sie in 3\,Sätzen. Welche \emph{Annahme} steckt in Ihrer Entscheidung?
\end{enumerate}

\ytquery{Plattform KPI Optimierung Conversion Checkout Marktplatz Steuerung}

\antwortlinien{20}

\clearpage

% --- AUFGABE 6 -----------------------------------------------------------
\aufgabenkopf{6}{Sechs Erfolgsfaktoren \textendash{} Selbst-Check}{\nivLii}{20\,min}

\begin{infobox}[Sechs Erfolgsfaktoren der Plattformimplementierung]
Aus dem Lehrskript Tag 2:
\begin{itemize}
  \item \textbf{Klare Zielsetzung} (welches Kundenproblem, welche Rolle im Vertrieb?)
  \item \textbf{Datenqualität} (Produktdaten, Kundendaten, CRM)
  \item \textbf{Kundennutzen} (Bedienbarkeit, Verfügbarkeit, Service)
  \item \textbf{Technische Stabilität} (Checkout, Schnittstellen, Performance)
  \item \textbf{Prozessklarheit} (Verantwortlichkeiten, Übergaben)
  \item \textbf{Akzeptanz im Vertrieb} (kein Schattenwettbewerb zum Außendienst)
\end{itemize}
\end{infobox}

\noindent\textbf{Arbeitsauftrag:}

\noindent Stellen Sie sich vor, Sie sind interne:r Berater:in bei SmartOffice24.

\begin{enumerate}
  \item Bewerten Sie jeden der sechs Erfolgsfaktoren mit einer Schulnote (1 = stark, 6 = ungenügend) basierend auf den Informationen aus Aufgabe 4/5. Wo der Fall keine Information enthält: machen Sie eine \emph{begründete} Annahme.
  \item Identifizieren Sie die \emph{zwei schwächsten} Faktoren \textendash{} hier liegt der größte Hebel.
  \item Entwickeln Sie für diese zwei Faktoren je \textbf{eine konkrete Maßnahme} in den nächsten 90\,Tagen.
  \item Schreiben Sie einen Schlusssatz: Welcher Erfolgsfaktor wird in Mittelstandsprojekten \emph{systematisch} unterschätzt \textendash{} und warum?
\end{enumerate}

\ytquery{Plattform Erfolgsfaktoren Datenqualität Vertriebsakzeptanz B2B Implementierung}

\antwortlinien{16}

\clearpage

% --- AUFGABE 7 -----------------------------------------------------------
\aufgabenkopf{7}{BauTech Connect \textendash{} drei Kanäle, ein Engpass}{\nivLii}{30\,min}

\begin{infobox}[Fallstudie: BauTech Connect]
``BauTech Connect'' ist ein mittelständischer Anbieter technischer Komponenten für Handwerksbetriebe, Bauunternehmen und Facility-Management-Dienstleister. Vor sechs Monaten wurde eine Plattformstrategie aus drei Bausteinen gestartet: \emph{eigener B2B-Shop}, \emph{zwei externe Marktplatzlistungen}, \emph{erste digitale Marketingkampagnen}.

\smallskip
\textbf{Aktuelle Kennzahlen (6 Monate):}
\begin{itemize}
  \item \emph{Eigener B2B-Shop:} 50.000 Besucher \textperiodcentered{} 2.200 registrierte Kunden \textperiodcentered{} 620 Bestellungen \textperiodcentered{} Warenkorb 760\,\euro{} \textperiodcentered{} CR 1{,}24\,\% \textperiodcentered{} Wiederkaufrate 18\,\% \textperiodcentered{} hoher Anteil erklärungsbedürftiger Produkte.
  \item \emph{Externe Marktplätze:} 180.000 Produktaufrufe \textperiodcentered{} 1.900 Bestellungen \textperiodcentered{} Warenkorb 230\,\euro{} \textperiodcentered{} starker Preisvergleich \textperiodcentered{} kaum Datenzugang \textperiodcentered{} niedrigere Marge.
  \item \emph{Digitale Kampagnen:} hoher Traffic, viele unqualifizierte Leads, Vertrieb beklagt geringe Abschlusswahrscheinlichkeit.
\end{itemize}

\smallskip
\textbf{Rahmenbedingungen:} 32\,Mio.\,\euro{} Jahresumsatz \textperiodcentered{} 1,4\,Mio.\,\euro{} Plattformumsatz bisher \textperiodcentered{} 300.000\,\euro{} Optimierungsbudget \textperiodcentered{} 9\,Monate \textperiodcentered{} Produktdatenqualität uneinheitlich \textperiodcentered{} CRM nur teilweise angebunden.
\end{infobox}

\noindent\textbf{Arbeitsauftrag:}

\begin{enumerate}
  \item Bewerten Sie die drei Bereiche (\emph{Shop}, \emph{Marktplätze}, \emph{Kampagnen}) anhand von \textbf{vier Kriterien}: Wachstum, Marge, strategische Kontrolle, Datenqualität. Skala 1\,--\,5.
  \item Identifizieren Sie die drei \textbf{größten Engpässe} in der Implementierung. Wo liegen sie wirklich \textendash{} bei Technik, Daten, Prozessen oder Rollen?
  \item Schlagen Sie eine \textbf{priorisierte Optimierungsroadmap} für 9\,Monate vor (3 Phasen je 3\,Monate).
  \item Begründen Sie, welche Aktivität bewusst \emph{nicht} ausgebaut wird.
\end{enumerate}

\ytquery{BauTech B2B Shop Marktplatz Kampagne Plattform Engpass Optimierung Roadmap}

\antwortlinien{22}

\clearpage

% --- AUFGABE 8 -----------------------------------------------------------
\aufgabenkopf{8}{Vier Zukunftstrends \textendash{} konkret übersetzen}{\nivLii}{20\,min}

\begin{infobox}[Vier Zukunftstrends im Plattformvertrieb]
\begin{enumerate}
  \item KI-gestützte Plattformsteuerung und automatisierte Vertriebsprozesse
  \item Personalisierung von Angeboten, Preisen und Kundenerlebnissen
  \item Datenbasierte Entscheidungsmodelle (statt Bauchgefühl) im digitalen Vertrieb
  \item Stärkere Bedeutung von Ökosystemen, Partnernetzwerken und Plattform-Governance
\end{enumerate}
\end{infobox}

\noindent\textbf{Arbeitsauftrag:}

\noindent Übersetzen Sie diese vier Trends von der ``Modeschlagworte''-Ebene auf die Ebene konkreter Vertriebsentscheidungen.

\begin{enumerate}
  \item Wählen Sie für jeden Trend \emph{einen} konkreten Anwendungsfall, der innerhalb von 12 Monaten bei BauTech Connect umsetzbar wäre.
  \item Bewerten Sie für jeden Anwendungsfall: \emph{Wer entscheidet?} \textendash{} \emph{Was kostet es ungefähr?} \textendash{} \emph{Wann wird Wirkung sichtbar?}
  \item Welcher Trend ist aus Ihrer Sicht \emph{überschätzt}, welcher \emph{unterschätzt}? Begründen Sie kurz.
  \item Formulieren Sie eine Faustregel (1\,--\,2 Sätze): Wann sollte ein Mittelständler einen Plattformtrend aktiv adoptieren \textendash{} und wann lieber abwarten?
\end{enumerate}

\ytquery{KI Plattform Vertrieb Personalisierung Datenmodell Ökosystem Trend B2B}

\antwortlinien{16}

\clearpage

% =========================================================================
% L3 AUFGABEN
% =========================================================================
\section*{Niveau L3 \textendash{} Management \& Entscheidungsarchitektur}
\addcontentsline{toc}{section}{L3 -- Management}

% --- AUFGABE 9 -----------------------------------------------------------
\aufgabenkopf{9}{BauTech Connect \textendash{} 9-Monats-Optimierungsentscheidung}{\nivLiii}{45\,min}

\begin{infobox}[Fallstudie: BauTech Connect (Fortsetzung)]
Ausgangslage wie in Aufgabe 7. Zusätzlich:
\begin{itemize}
  \item Die Geschäftsführung erwartet \textbf{eine klare Empfehlung} zur Weiterentwicklung der Plattformstrategie.
  \item Das Optimierungsbudget von \textbf{300.000\,\euro} ist für die nächsten 9 Monate verfügbar.
  \item Vertrieb, Marketing und IT bewerten den Erfolg unterschiedlich \textendash{} es gibt keine gemeinsamen KPIs.
  \item Kundenservice meldet zunehmend Rückfragen zu digitalen Bestellungen.
  \item Der Vertrieb befürchtet, dass Bestandskunden zunehmend nur noch Preise vergleichen.
\end{itemize}

\smallskip
\textbf{Verfügbare Hebel:}
\begin{itemize}
  \item Produktdaten und Shop-Optimierung
  \item CRM- und Datenarchitektur (Anbindung, Lead Scoring)
  \item Kampagnen (Zielsegmente, Performance, Bewertung)
  \item Marktplatzsteuerung (Sortiment, Marge, Steuerung)
  \item Personalisierung, Wiederkauflogik, Management-Dashboard
\end{itemize}
\end{infobox}

\noindent\textbf{Arbeitsauftrag:}

\begin{enumerate}
  \item Treffen Sie als \emph{Chief Digital Sales Officer} eine \textbf{priorisierte Optimierungsentscheidung} für die nächsten 9 Monate. Strukturieren Sie sie in drei Phasen (Monat 1\,--\,3, 4\,--\,6, 7\,--\,9).
  \item Verteilen Sie das Budget von \textbf{300.000\,\euro} auf maximal fünf Hebel. Begründen Sie jede Budgetlinie.
  \item Wählen Sie \textbf{einen Zukunftstrend} aus Aufgabe 8 und erklären Sie, wie er in die Roadmap eingebaut wird (oder \emph{warum} er bewusst herausgehalten wird).
  \item Beschreiben Sie, wie sich die \textbf{Rollen} von Außendienst, Innendienst, Marketing und Kundenservice im Lauf der 9 Monate verändern.
  \item Treffen Sie \emph{mindestens eine Entscheidung unter Unsicherheit} \textendash{} und benennen Sie ausdrücklich die Annahme, die dafür stimmen muss.
  \item Definieren Sie \textbf{drei Management-KPIs}, die nach 9 Monaten zwingend im Geschäftsführungs-Meeting berichtet werden sollen.
\end{enumerate}

\ytquery{Plattform Optimierungsroadmap Budget Mittelstand CDSO Entscheidung 9 Monate}

\antwortlinien{32}

\clearpage

% --- AUFGABE 10 ----------------------------------------------------------
\aufgabenkopf{10}{Zielkonflikt \textendash{} Reichweite versus Plattformkontrolle}{\nivLiii}{30\,min}

\noindent\textbf{Diskussionsaufgabe.}

\noindent Im BauTech-Fall (Aufgabe 7/9) liegt ein klassischer Zielkonflikt offen: Die externen Marktplätze liefern \emph{Volumen} (1.900 Bestellungen in 6 Monaten), schwächen aber \emph{Marge}, \emph{Kundendatenzugang} und \emph{Differenzierung}. Der eigene B2B-Shop liefert geringere Reichweite, aber höheren Warenkorb, bessere Daten und größere strategische Kontrolle.

\medskip
\noindent\textbf{Arbeitsauftrag:}

\begin{enumerate}
  \item Beschreiben Sie in eigenen Worten den Zielkonflikt zwischen \emph{kurzfristiger Reichweite} und \emph{langfristiger Plattformkontrolle} (3\,--\,4 Sätze).
  \item Entwickeln Sie eine \textbf{Entscheidungsregel}, wann ein Mittelständler aktiv auf einem externen Marktplatz präsent sein \emph{soll} \textendash{} und wann er ihn verlassen (oder gar nicht erst betreten) sollte. Nutzen Sie dafür mindestens drei Kriterien.
  \item Wählen Sie \emph{ein konkretes Sortimentssegment} bei BauTech (z.\,B.\ Standard-Verbindungsmittel, Premium-Werkzeuge, Spezialkomponenten) und entscheiden Sie für dieses Segment: Marktplatz \emph{ja} oder \emph{nein}? Begründen Sie.
  \item Formulieren Sie eine Faustregel (1\,--\,2 Sätze): Welche Mindestbedingung muss ein eigener Direktkanal erfüllen, bevor man externe Marktplätze überhaupt aufbaut?
\end{enumerate}

\ytquery{Marktplatz Eigenkanal Reichweite Marge Plattformkontrolle B2B Entscheidung}

\antwortlinien{20}

\clearpage

% --- AUFGABE 11 ----------------------------------------------------------
\aufgabenkopf{11}{Transferprojekt \textendash{} Plattformarchitektur 2030}{\nivLiii}{45\,min}

\begin{infobox}[Rolle und Ausgangslage]
\textbf{Ihre Rolle:} Chief Digital Sales Officer (oder Plattformstrategie-Architekt:in) eines mittelständischen B2B-Unternehmens.

\smallskip
\textbf{Ausgangslage:} Erste Plattformaktivitäten sind aufgebaut: eigener Webshop, externe Marktplatzlistungen, digitale Kampagnen, einfaches CRM. Es gibt Wachstum \emph{und} Probleme: steigender Preisdruck, unklare Datenverantwortung, parallele Prozesse, zunehmende Abhängigkeit von externen Plattformen.

\smallskip
Die Geschäftsführung fordert eine Plattformstrategie 2030, die Wachstum ermöglicht und gleichzeitig Kundenschnittstelle, Datenhoheit, Marge und Markenpositionierung sichert.

\smallskip
\textbf{Rahmenbedingungen:}
\begin{itemize}
  \item Budget begrenzt.
  \item Digitale Wettbewerber wachsen schneller.
  \item Interne Datenqualität unzureichend.
  \item Vertrieb, Marketing und IT verfolgen teilweise unterschiedliche Ziele.
  \item Kundenerwartungen: schneller, personalisierter, digitaler.
  \item KI, Automatisierung und Plattformökosysteme entwickeln sich dynamisch.
  \item Regulatorische Anforderungen an Daten und faire Plattformpraktiken nehmen zu.
\end{itemize}
\end{infobox}

\noindent\textbf{Arbeitsauftrag:} Entwickeln Sie eine strategische \emph{Entscheidungsarchitektur} \textendash{} keine reine Maßnahmenliste \textendash{} für die Implementierung, Optimierung und Weiterentwicklung der Plattformstrategie bis 2030. Erarbeiten Sie schriftlich:

\begin{enumerate}
  \item Ein \textbf{Zielbild für 2030} (4\,--\,5 Sätze).
  \item Eine \textbf{Bewertung der bestehenden Plattformkanäle} nach Wachstum, Marge, Datenzugang, Kundenbindung und strategischer Kontrolle.
  \item Eine \textbf{Entscheidung}, welche Plattformaktivitäten ausgebaut, optimiert, begrenzt oder beendet werden.
  \item Eine \textbf{KPI-Architektur} für die Managementsteuerung (operative KPIs vs.\ Management-KPIs).
  \item Eine \textbf{Governance-Struktur} mit klaren Verantwortlichkeiten zwischen Vertrieb, Marketing, IT, Datenmanagement und Geschäftsführung.
  \item Eine \textbf{Roadmap} (12 / 24 / 36+\,Monate).
  \item Eine \textbf{Einschätzung relevanter Zukunftstrends} und deren strategischer Bedeutung.
  \item Eine \textbf{Risikoarchitektur} zu Plattformabhängigkeit, Datenverlust, Preisdruck, technischer Komplexität, organisatorischer Überforderung.
  \item \textbf{Drei Managemententscheidungen unter Unsicherheit} \textendash{} jeweils mit der Annahme, die stimmen muss.
\end{enumerate}

\noindent\textbf{Zusatzanforderung:} Treffen Sie mindestens \emph{eine Entscheidung, bei der Sie kurzfristiges Wachstum bewusst zugunsten langfristiger strategischer Kontrolle begrenzen}. Begründen Sie aus Senior-Level-Perspektive.

\ytquery{Plattformstrategie 2030 Mittelstand Architektur Governance KPI Zukunftstrends}

\antwortlinien{38}

\clearpage

% --- Abschluss -----------------------------------------------------------
\section*{Abschluss}
\thispagestyle{plain}

\noindent Sie haben alle elf Aufgaben bearbeitet. Vergleichen Sie nun Ihre Antworten mit dem \emph{Lösungsheft Tag 2}.

\medskip
\begin{infobox}[Was Sie jetzt können sollten]
\begin{itemize}
  \item Plattformauswahl, -implementierung und -optimierung sauber abgrenzen.
  \item Die sechs Phasen einer Plattformimplementierung benennen und auf SmartOffice24 anwenden.
  \item Plattform-KPIs nach Sichtbarkeit, Conversion und Wirtschaftlichkeit unterscheiden.
  \item Zielkonflikte zwischen Reichweite, Marge, Kontrolle und Datenzugang bewerten.
  \item Zukunftstrends (KI, Personalisierung, Ökosysteme) in konkrete Vertriebsentscheidungen übersetzen.
  \item Eine 9-Monats-Optimierungsroadmap mit Budgetlogik formulieren.
  \item Eine Plattformarchitektur 2030 als Entscheidungsarchitektur denken \textendash{} nicht als Maßnahmenliste.
\end{itemize}
\end{infobox}

\vspace{1cm}
\noindent{\color{lpAccent}\rule{\linewidth}{0.6pt}}\\[4pt]
{\footnotesize\sffamily\color{lpGray}Lernplatt \textperiodcentered{} by Can Siebert \textperiodcentered{} Kurs 22 (Bo 143) \textperiodcentered{} Tag 2 \textperiodcentered{} Aufgabenheft \textperiodcentered{} \textcopyright{} 2026}

\end{document}
